\documentclass[10pt,a4paper]{amsart}
\usepackage[margin=19mm]{geometry}
\usepackage{amsmath,amssymb,mathtools}
\usepackage[scr=rsfs,cal=euler]{mathalfa}
\usepackage{xfrac}
\usepackage[unicode,hidelinks,bookmarks=false]{hyperref}
\newcommand{\ie}{i.e.\ }
\newcommand{\cf}{cf.\ }
\newcommand{\resp}{respectively\ }
\newcommand{\Subsection}{\subsection}
\providecommand{\nobreakdash}{\nobreak\hbox{-}}
\setlength{\emergencystretch}{2em}
\multlinegap0pt
\usepackage{bm}
\usepackage{bbm}
\usepackage{dsfont}
\usepackage{xspace}
\usepackage{cancel}
\usepackage{mathtools}
\usepackage{amsthm}
\makeatletter
\@ifundefined{theorem}{\newtheorem{theorem}{Theorem}}{}
\@ifundefined{lemma}{\newtheorem{lemma}[theorem]{Lemma}}{}
\@ifundefined{proposition}{\newtheorem{proposition}[theorem]{Proposition}}{}
\makeatother
\newcommand{\LT}{\operatorname{LT}}
\title[Reduction of the Finite Generation Problem for Leading Term ]{Reduction of the Finite Generation Problem for Leading Term Ideals under Arbitrary Rational Monomial Orders to the Lexicographic}
\author{Xiaopeng Zheng}
\date{Source: arXiv:2607.25372v1 (CC BY 4.0)}
\begin{document}
\maketitle

\noindent\emph{Source and licence.} Reduction of the Finite Generation Problem for Leading Term Ideals under Arbitrary Rational Monomial Orders to the Lexicographic. Version arXiv:2607.25372v1 (CC BY 4.0), distributed under the Creative Commons Attribution 4.0 International licence, as recorded in the arXiv licence metadata for this exact version and shown on the abstract page https://arxiv.org/abs/2607.25372v1. Full rights notice: licenses/arxiv-2607.25372-CC-BY-4.0.txt.\par\smallskip

\noindent\textit{Editorial scope.} Counted unit: the Theorem whose source label is \texttt{thm:fixed-ideal-reduction} in the article (source0.tex lines 859--867, source label \texttt{thm:fixed-ideal-reduction}), delivered together with the article's own complete original proof of it (source0.tex lines 869--912, one \texttt{proof} environment of 44 lines). No mathematics is written, repaired or re-derived: statement and proof are the article's own text line for line. Required context retained uncounted: prop:finite-grobner-basis (lines 303--308); prop:homogeneous-lifts (lines 772--783); thm:exact-leading-term-correspondence (lines 812--822). Every definition, display and label of those lines is retained unchanged. Omitted from this package and not redistributed: 1--302, 309--771, 784--811, 823--858, 868--868, 913--1120. Nothing inside a retained excerpt is omitted or altered: every retained body is delivered exactly as the article's lines are written, so no omission receipt of any kind accompanies this package. Citations: the retained excerpts carry no citation command at all, so no bibliography entry of the article is transcribed and no reference is redistributed. Reference adaptation: none was needed and none was made; every reference inside the delivered unit resolves inside it, so no unresolved reference or citation is produced. Printed slips kept literal: no printed word, symbol, hypothesis or number of the counted statement or of its proof was altered. Layout and class: the article's own class is \texttt{elsarticle} with packages including amsmath, amssymb, amsthm, mathtools, hyperref; the wrapper is the lane's pinned \texttt{amsart} editorial wrapper at 10pt on A4, which restates the theorem-like environments the retained text needs and adds the article's own macro definitions that the retained text needs (\texttt{\textbackslash LT}), with extra wrapper packages (bm, bbm, dsfont, xspace, cancel, mathtools). The delivered numbering is the unit's own. Assets: no retained span contains a figure environment, a graphics-inclusion command, a PDF-page inclusion, a table or a TikZ picture; no image file is copied into this package, the package declares no asset dependency and no image file is redistributed with it. Licence: version arXiv:2607.25372v1 of this article is distributed under the Creative Commons Attribution 4.0 International licence, as recorded in the arXiv licence metadata for this exact version and shown on the abstract page https://arxiv.org/abs/2607.25372v1. A full rights notice is delivered at licenses/arxiv-2607.25372-CC-BY-4.0.txt. Characters: every retained excerpt is pure ASCII, so the delivered engine, XeLaTeX with its default Latin Modern text font, reports no missing character.
	\begin{proposition}[Finite Gr\"obner bases and leading term ideals]
		\label{prop:finite-grobner-basis}
		Let $I\subseteq A$ be a finitely generated ideal. Then
		$\LT_\prec(I)$ is finitely generated if and only if $I$ admits a
		finite Gr\"obner basis with respect to $\prec$.
	\end{proposition}

	\begin{proposition}[Homogeneous lifts]
		\label{prop:homogeneous-lifts}
		With the notation above, the following statements hold.
		\begin{enumerate}
			\item For every $f\in A$, the polynomial $\Psi_M(f)$ is homogeneous
			of degree zero.
			\item The ideal $H$ is graded.
			\item $\chi(H)=I$.
			\item Every $g\in I$ has the homogeneous lift $\Psi_M(g)\in H$, and
			$\chi(\Psi_M(g))=g$.
		\end{enumerate}
	\end{proposition}

	\begin{lemma}[Exact dehomogenization of leading terms]
		\label{thm:exact-leading-term-correspondence}
		Let $0\neq P\in S$ be homogeneous for the grading $\deg_M$. Then
		$\chi(P)\neq0$ and
		\begin{equation}
			\label{eq:exact-leading-term-correspondence}
			\chi\!\left(\LT_{\prec_{\mathrm{lex}}}(P)\right)
			=
			\LT_\prec\!\left(\chi(P)\right).
		\end{equation}
	\end{lemma}

	\begin{theorem}[Reduction to the lexicographic order]
		\label{thm:fixed-ideal-reduction}
		Let $I=\langle f_1,\ldots,f_s\rangle\subseteq A$, and let $H\subseteq
		S$ be the ideal constructed above. If
		$\LT_{\prec_{\mathrm{lex}}}(H)$ is finitely generated, then
		$\LT_\prec(I)$ is finitely generated. Equivalently, if $H$ admits a
		finite Gr\"obner basis with respect to $\prec_{\mathrm{lex}}$, then
		$I$ admits a finite Gr\"obner basis with respect to $\prec$.
	\end{theorem}

	\begin{proof}
		The zero ideal is immediate, so assume $I\neq0$. Since $H$ is
		finitely generated, Proposition~\ref{prop:finite-grobner-basis}
		and the hypothesis yield a finite Gr\"obner basis
		$Q_1,\ldots,Q_t$ of $H$. For each $j$, let $P_j$ be
		the homogeneous component of $Q_j$ that contains
		$\LT_{\prec_{\mathrm{lex}}}(Q_j)$. Since $H$ is graded, one has
		$P_j\in H$ and
		$\LT_{\prec_{\mathrm{lex}}}(P_j)=
		\LT_{\prec_{\mathrm{lex}}}(Q_j)$. Thus the leading terms of
		$P_1,\ldots,P_t$ still generate
		$\LT_{\prec_{\mathrm{lex}}}(H)$.
		
		Set $p_j=\chi(P_j)$. Then $p_j\in I$, and
		Lemma~\ref{thm:exact-leading-term-correspondence} gives
		$p_j\neq0$ and
		$\LT_\prec(p_j)=
		\chi(\LT_{\prec_{\mathrm{lex}}}(P_j))$. Let $0\neq g\in I$. Proposition~\ref{prop:homogeneous-lifts} shows
		that $\Psi_M(g)$ is a nonzero homogeneous element of $H$. Since the
		leading terms of the $P_j$ generate
		$\LT_{\prec_{\mathrm{lex}}}(H)$, there exist
		$A_1,\ldots,A_t\in S$ such that
		\[
		\LT_{\prec_{\mathrm{lex}}}(\Psi_M(g))
		=
		\sum_{j=1}^t A_j\LT_{\prec_{\mathrm{lex}}}(P_j).
		\]
		Applying $\chi$ and using
		Lemma~\ref{thm:exact-leading-term-correspondence} gives
		$\LT_\prec(g)=\sum_{j=1}^t\chi(A_j)\LT_\prec(p_j)$.
		Therefore $\LT_\prec(I)\subseteq
		\langle\LT_\prec(p_1),\ldots,\LT_\prec(p_t)\rangle$. The reverse
		inclusion follows from $p_j\in I$, and hence
		\begin{equation}
			\label{eq:leading-term-ideal-descent}
			\LT_\prec(I)
			=
			\left\langle\LT_\prec(p_1),\ldots,\LT_\prec(p_t)\right\rangle.
		\end{equation}
		Finally, the finite set
		$\{f_1,\ldots,f_s,p_1,\ldots,p_t\}$ generates $I$, and its leading
		terms generate $\LT_\prec(I)$. It is therefore a finite Gr\"obner
		basis of $I$.
	\end{proof}

\end{document}
